% !TeX root = ../main.tex
\section{模型评价、改进与推广}

\subsection{模型优点}

第一，四问采用递进且可精确求解的模型链：问题一用离散精确尾概率控制两类误判风险；问题二以联合信念状态 MDP 处理信息状态相关决策；问题三把大规模决策明确限制为固定策略类，以完全枚举获得固定策略空间内的全局最优，并用谱半径判据刻画可行性；问题四以 Beta 后验情景平均把抽样不确定性纳入决策。各问结果均可复现，数值验证与参数敏感性相互衔接。

第二，问题二与问题三的模型动机清晰：问题二证明“不合格成品”观测会使两部件后验相关、固定 0-1 决策无法表达该信息反馈；问题三在状态空间随工序增长时主动限制策略类，并解释 74\% 策略不可行的“缺陷回流闭环”机制，使模型选择与生产机制对应。

第三，问题四把“抽样只提供有限证据”以标准贝叶斯框架显式建模：后验期望利润作为风险中性决策准则，标准差与 $5\%$ 分位数辅助刻画波动与下行风险；样本量增大时结果收敛到点估计口径，先验与情景数稳定性检验支持主要结论。

\subsection{模型局限}

问题一未获得总体量 $N$、不可接受次品率与目标检验功效，样本量口径是题目信息不足下的建模选择；检测误差与序贯抽样未纳入。问题二的信念状态数量随部件数增长较快；问题三的固定策略类不随信息状态改变，其敏感性为五类参数的单因素 $\pm20\%$ 网格，联合扰动未覆盖。问题四的实际抽样样本量未知，$n=40$ 为代表性情景而非题面事实；在均匀先验下情形 5 的策略切换对先验选择敏感，且完整信念状态 MDP 的参数不确定性版本因规模限制未求解。若需要风险厌恶决策，可将目标改为后验期望利润减风险惩罚项或 CVaR 准则。

\subsection{模型推广}

问题一的“离散尾概率—三态判定”可用于其他批次验收问题，替换标称质量水平、风险上限与抽样机制即可。问题二的“联合信念—贝叶斯更新—Bellman 决策”可扩展到更多零配件或多级半成品，状态规模可通过因子化信念、状态聚合或近似动态规划控制。问题三的“期望递推—精确枚举—启发式扩展”可推广到一般树状或网状装配结构，递推按新结构的输入输出重写，遗传算法编码按决策位数扩展。问题四的后验情景平均框架可替换先验、样本量或多参数联合后验，把待估计参数替换为先验分布后即可复用内层解析期望评估器。
